\section{物理量、函数与图像}
\subsection{先明确哪个量依赖哪个量}
把位置写成 $x(t)$，表示给定时间 $t$ 后可以确定位置；这里 $t$ 是自变量，$x$ 是因变量。函数关系不一定来自一个已经解出的显式公式，也可以来自几何约束、实验测量或动力学方程。开始计算前，应交代自变量范围、物理量正方向和参数的含义。例如绳长 $L$ 在一次运动中可视为固定参数，时间在变；若题目改为比较不同绳长的装置，$L$ 才成为另一个需要改变的量。

定义域是物理模型的一部分。小球尚未碰地时的自由落体公式不能直接延伸到碰地之后；液面高度必须非负；某个平方根的两个代数分支也未必都对应实际装配位置。几何图帮助确定允许范围，随后求导、积分和取极值都应在同一范围内进行。

函数图的横纵轴可以是不同物理量，因此纸面上的一个角度不直接等于具有量纲的变化率。设运动关系为 $x=L(t/\tau)^2$，其中 $L$ 为长度尺度、$\tau$ 为时间尺度。令 $s=t/\tau$，就有 $x/L=s^2$。图\ref{c1foundation:scales}左边使用物理坐标，右边使用无量纲坐标；二者表示同一过程，但横纵轴的刻度不同。
\fig{c1foundation:scales}{物理坐标与无量纲坐标表示同一变化关系}{\begin{tikzpicture}[scale=.8]
\begin{scope}
\draw[axis](0,0)--(3,0)node[right]{$t$};\draw[axis](0,0)--(0,3)node[above]{$x$};
\draw[world,domain=0:2.4,samples=50]plot(\x,{.45*\x*\x});
\draw[guide](1.8,0)node[below]{$\tau$}--(1.8,1.458)--(0,1.458)node[left]{$L$};
\node at(1.4,-.9){$x=L(t/\tau)^2$};
\end{scope}
\begin{scope}[xshift=5.7cm]
\draw[axis](0,0)--(3,0)node[right]{$t/\tau$};\draw[axis](0,0)--(0,3)node[above]{$x/L$};
\draw[world,domain=0:1.6,samples=50]plot(\x,{\x*\x});
\draw[guide](1,0)node[below]{$1$}--(1,1)--(0,1)node[left]{$1$};
\node at(1.4,-.9){同一过程的尺度表示};
\end{scope}
\end{tikzpicture}}

求速度时，$L,\tau$保持不变，而$s=t/\tau$随$t$变化。沿图\ref{c1foundation:scales}两套坐标之间的对应关系，链式法则给出
\[
v=\frac{\dd x}{\dd t}
=L\frac{\dd(s^2)}{\dd s}\frac{\dd s}{\dd t}
=L(2s)\frac1\tau
=\frac{2Lt}{\tau^2}.
\]
无量纲图上的导数$2s$不能直接当作有速度单位的$v$；必须乘回$L/\tau$。这一点会在坐标变换、实验图像斜率和近似展开中反复使用。

\subsection{平均变化与瞬时变化回答不同问题}
在有限时间$\Delta t\ne0$内，平均速度是$[x(t+\Delta t)-x(t)]/\Delta t$。对上面的二次位置函数，先展开有限增量：
\begin{align*}
\frac{x(t+\Delta t)-x(t)}{\Delta t}
&=\frac{L}{\tau^2}\frac{(t+\Delta t)^2-t^2}{\Delta t}\\
&=\frac{L}{\tau^2}\frac{2t\Delta t+(\Delta t)^2}{\Delta t}
=\frac{2Lt}{\tau^2}+\frac{L\Delta t}{\tau^2}.
\end{align*}
第二项来自有限时间内速度继续变化；只有令时间间隔趋零，才得到时刻$t$的瞬时速度。这个过程不是把分母直接写成零，而是先在非零间隔下整理，再取极限。

反过来，已知每个时刻的速度，要恢复一段过程的位移，需要累积 $v\,\dd t$，还要给出初始位置。导数、积分和初始条件各自提供不同信息：导数描述局部变化率，积分连接有限过程，初值确定常数。后文的严格定义将把这些图像变成可检验的计算，而不改变这里每个物理量与图中轴的对应。
